\section{Células diádicas e operador gerador}
\subsection{Células, posse e centros}
Para $k\in\N_0$ e $m\in\Z$, sejam
\eqn{09}{J_{k,m}=[m2^{-k},(m+1)2^{-k}],\qquad J_{k,m}^{\rm own}=[m2^{-k},(m+1)2^{-k}),}
e
\eqn{10}{x_0(k,m)=\left(m+\frac12\right)2^{-k}=\frac{2m+1}{2^{k+1}}.}
A célula fechada sustenta o segmento; a versão semiaberta atribui cada fronteira a um único proprietário. Escreva
\[
\D=\{j2^{-K}:j\in\Z,\ K\in\N_0\},\qquad \D_{\rm mid}=\D\setminus\Z.
\]
\begin{lema}[estrutura diádica]
Para cada $k$, as células $J_{k,m}^{\rm own}$ particionam $\R$ e
\eqn{11}{J_{k,m}=J_{k+1,2m}\cup J_{k+1,2m+1},\qquad |J_{k,m}|=2^{-k}.}
O mapa $(k,m)\mapsto x_0(k,m)$ é uma bijeção de $\N_0\times\Z$ sobre $\D_{\rm mid}$, que é denso em $\R$.
\end{lema}
\begin{proof}
Para $x\in\R$, o único índice de posse na escala $k$ é $m=\lfloor2^kx\rfloor$; isso prova a partição. A identidade de refinamento decorre de cindir $[2m2^{-(k+1)},(2m+2)2^{-(k+1)}]$ em seu ponto médio.

Pela \eqref{eq:10}, o centro tem numerador ímpar e denominador reduzido $2^{k+1}\ge2$, logo pertence a $\D_{\rm mid}$. Reciprocamente, cada $x\in\D_{\rm mid}$ possui forma reduzida única $j/2^K$, com $K\ge1$ e $j$ ímpar; tomando $k=K-1$ e $m=(j-1)/2$, obtém-se $x=x_0(k,m)$. A mesma unicidade prova a injetividade. Por fim, os centros de escala $k$ estão espaçados por $2^{-k}$; escolhendo o mais próximo de um real dado, a distância tende a zero com $k$, o que prova a densidade.
\end{proof}
\figura{02_celulas}{.92}{Células diádicas, centros e posse das fronteiras. Cada célula se divide em dois filhos; os centros determinam a célula de origem.}{celulas}

\subsection{Retas ancoradas e descompactação}
Dados $x_0,y_0\in\R$ e $s\ne0$, a reta ancorada é
\eqn{12}{\ell(x)=y_0+s(x-x_0)=a+bx,\qquad b=s,\quad a=y_0-sx_0.}
Para $n\in\N_0$, aplique sucessivamente a inversa do pareamento:
\eqn{13}{(u,v)=\pi^{-1}(n),\qquad(k,\mu)=\pi^{-1}(u),\qquad(r,t)=\pi^{-1}(v).}
\begin{lema}
O mapa $n\mapsto(k,\mu,r,t)$ de \eqref{eq:13} é uma bijeção de $\N_0$ sobre $\N_0^4$.
\end{lema}
\begin{proof}
Sua inversa explícita é
\[
(k,\mu,r,t)\longmapsto\pi\bigl(\pi(k,\mu),\pi(r,t)\bigr),
\]
uma composição de bijeções pelo lema~2.1.
\end{proof}
Defina
\eqn{14}{m=Z(\mu),\quad x_{0,n}=x_0(k,m),\quad y_{0,n}=Q(r),\quad s_n=Q^*(t),}
e converta para a forma global por
\eqn{15}{b_n=s_n,\qquad a_n=y_{0,n}-s_nx_{0,n}.}
Se $x_L=m2^{-k}$ e $x_R=(m+1)2^{-k}$, ponha
\eqn{16}{c_n=\min\{\ell_n(x_L),\ell_n(x_R)\},\qquad d_n=\max\{\ell_n(x_L),\ell_n(x_R)\}.}
\begin{definicao}[gerador e biblioteca]
Se $\mathcal I$ é o conjunto dos intervalos fechados não degenerados de $\R$, defina
\eqn{17}{\Gamma:\N_0\to\R\times(\R\setminus\{0\})\times\mathcal I,\qquad
\Gamma(n)=(a_n,b_n,[c_n,d_n]).}
A biblioteca é $\LL=\Gamma(\N_0)$. O átomo de endereço $n$ é a restrição de $\ell_n$ à célula $J_{k(n),m(n)}$.
\end{definicao}
Como $s_n\ne0$ e $x_L\ne x_R$, o intervalo em \eqref{eq:16} é não degenerado. A inclinação não nula torna a interface invertível no eixo espacial: uma peça constante teria imagem pontual e perderia a informação da célula. Para $\Gamma(n)=(a_n,b_n,[c_n,d_n])$, a pré-imagem espacial é
\eqn{18}{I_n^x=\left[
\min\left\{\frac{c_n-a_n}{b_n},\frac{d_n-a_n}{b_n}\right\},
\max\left\{\frac{c_n-a_n}{b_n},\frac{d_n-a_n}{b_n}\right\}\right].}

\subsection{Recuperação da célula e endereço canônico}
\begin{lema}[validade espacial exata]
Para todo $n\in\N_0$, o intervalo $I_n^x$ definido em \eqref{eq:18} satisfaz
\[
I_n^x=J_{k(n),m(n)}.
\]
\end{lema}
\begin{proof}
Como $b_n\ne0$, $\ell_n$ é uma bijeção afim estritamente monótona. Sua restrição a $[x_L,x_R]$ tem imagem $[c_n,d_n]$ por \eqref{eq:16}; logo $\ell_n^{-1}([c_n,d_n])=[x_L,x_R]$. Os dois quocientes em \eqref{eq:18} são precisamente $\ell_n^{-1}(c_n)$ e $\ell_n^{-1}(d_n)$, possivelmente em ordem invertida quando $b_n<0$.
\end{proof}
\begin{teorema}[unicidade do endereço]
$\Gamma$ é injetiva. Portanto, é uma bijeção de $\N_0$ sobre $\LL$ e cada átomo tem um único endereço.
\end{teorema}
\begin{proof}
Suponha $\Gamma(n)=\Gamma(n')=(a,b,[c,d])$. Pelo lema~3.4, a tripla recupera a célula comum $J=[x_L,x_R]$. Seu comprimento determina unicamente $k=-\log_2|J|$, e $m=2^kx_L$. Assim, também se recupera $x_0=(m+1/2)2^{-k}$. A inclinação local é $s=b$ e o valor ancorado é $y_0=a+bx_0$.

Pelas bijeções dos lemas~2.2, 2.4 e 3.1, os dados $(x_0,y_0,s)$ determinam unicamente $(k,\mu,r,t)$. O lema~3.2 determina então um único $n$. Logo $n=n'$.
\end{proof}
Em uma fronteira, duas células fechadas podem estar ativas. A célula semiaberta escolhe uma proprietária; em uma montagem contínua, as duas peças têm o mesmo valor no nó.
\begin{teorema}[dicionário de jatos racionais]
O mapa
\eqn{19}{n\longmapsto(x_{0,n},y_{0,n},s_n)}
é uma bijeção de $\N_0$ sobre $\D_{\rm mid}\times\Q\times\Q^*$. Sua imagem é densa em $\R^3$.
\end{teorema}
\begin{proof}
A bijetividade resulta da composição das bijeções em \eqref{eq:13}--\eqref{eq:14}. A densidade decorre da densidade de $\D_{\rm mid}$, de $\Q$ e de $\Q^*$ em $\R$, aplicada coordenada a coordenada na norma do máximo.
\end{proof}
O dicionário \eqref{eq:19} parametriza todos os 1-jatos com centro diádico não inteiro, valor racional e derivada racional não nula. Montagens contínuas ligam esses dados locais por valores de fronteira compartilhados.
